% =============================================================================
% ASP: The Alignment Security Protocol — The Synthetic Immunity Protocol
% IC3/Encode Club Shape Rotator Hackathon Submission
% =============================================================================

\documentclass[sigconf,natbib=false]{acmart}

\usepackage{algorithm}
\usepackage{algpseudocode}
\usepackage{amsmath}
\usepackage{amssymb}
\usepackage{booktabs}
\usepackage{listings}
\usepackage{xcolor}
\usepackage{hyperref}
\usepackage{microtype}
\usepackage{tikz}
\usetikzlibrary{arrows.meta, positioning, shapes.geometric}

% --- Custom algpseudocode blocks for try/finally ---
\algnewcommand\Try{\textbf{try:}}
\algnewcommand\Finally{\textbf{finally:}}
\algnewcommand\EndTry{}

% --- ACM metadata (required by acmart) ---
\acmConference[Shape Rotator '25]{IC3/Encode Club Hackathon}{2025}{Cornell University, Ithaca, NY}
\acmDOI{}
\acmISBN{}
\copyrightyear{2025}

% --- Code listing style ---
\lstset{
  basicstyle=\ttfamily\footnotesize,
  keywordstyle=\color{blue}\bfseries,
  commentstyle=\color{gray}\itshape,
  stringstyle=\color{teal},
  frame=single,
  breaklines=true,
  columns=fullflexible,
  keepspaces=true,
}

% =============================================================================
\begin{document}

\title{ASP: The Alignment Security Protocol\\
\large A Synthetic Immunity Architecture for Hardware-Enforced LLM Safety}

\author{Submitted to the IC3/Encode Club Shape Rotator Hackathon}
\affiliation{%
  \institution{Cornell Initiative for CryptoCurrencies \& Contracts}
  \city{Ithaca}
  \state{New York}
  \country{USA}
}
\email{anon@ic3.cornell.edu}

% =============================================================================
\begin{abstract}
Large language models (LLMs) deployed in production environments are subject
to adversarial attacks that systematically bypass safety alignment—including
adversarial roleplay, prompt/context injection, and system-prompt exfiltration.
Existing defenses (keyword filters, regex blocklists, fine-tuned refusal
classifiers) share a common failure mode: they operate on surface form, making
them vulnerable to paraphrase, encoding, and compositional obfuscation.

We present the \emph{Alignment Security Protocol} (ASP), a defense-in-depth
architecture that moves safety enforcement to three hardware-backed primitives:
(1) a Trusted Execution Environment (TEE) boundary that destructs raw prompts
after encoding, enforcing the invariant that only sanitized context ever reaches
the LLM; (2) a Morphological Intent Encoder that classifies threat level via
geometric proximity to attack manifolds in latent space, achieving surface-form
invariance; and (3) a threshold signature scheme (N-of-M consensus) that
prevents single-node compromise from forging a safe verdict.
An epidemic gossip protocol propagates new attack signatures as
\emph{vaccines} across the validator network in $O(\log N)$ rounds.

We prove that a single compromised validator cannot produce a
\textsc{Verified\_Immunity} verdict and that the encoder is invariant to
surface-form transformations that preserve semantic intent. ASP is
model-agnostic, implemented in Python~3.11 with no ML framework dependencies
in the core path, and exposes a JSON-RPC~2.0 telemetry interface for
interoperability.

\textbf{Keywords:} TEE isolation, threshold cryptography, federated AI safety,
morphological intent detection, jailbreak defense, gossip protocols,
Shamir secret sharing.
\end{abstract}

\maketitle

% =============================================================================
\section{Introduction}
\label{sec:intro}

The deployment of aligned LLMs at scale has produced an adversarial arms race.
Safety alignment, whether via RLHF~\cite{ouyang2022training},
Constitutional~AI~\cite{bai2022constitutional}, or supervised fine-tuning,
is applied at training time and encoded in the model's weight distribution.
At inference time, the model is a text transducer operating on a flat context
window, with no architectural mechanism to distinguish a system operator's
instruction from an adversarial user injection.

\subsection{The Alignment Tax Problem}
\label{sec:intro:tax}

Contemporary safety wrappers exhibit a trilemma we term the \emph{alignment
tax}: precision, recall, and latency cannot all be optimized simultaneously
against a capable adversary. Keyword and regex filters have near-zero latency
but suffer from recall failure under paraphrase, encoding substitution
(Base64, Unicode homoglyphs, ROT-$n$), and inter-language attacks. Neural
classifiers trained on known jailbreak datasets exhibit high recall on
in-distribution attacks but degrade sharply on compositional novelty—a
property formally demonstrated by Wei et al.~\cite{wei2023jailbroken} and
Zou et al.~\cite{zou2023universal}.

The root failure is architectural, not parametric: when a safety mechanism
operates on the same surface as the attack, the attacker's search space is
the same as the defender's. No filter defined over token sequences can be
closed under the equivalence class of paraphrases that preserve harmful
intent.

\subsection{Why Hardware-Enforced Enforcement}
\label{sec:intro:hw}

The NDAI Agreements framework~\cite{ndai2024tee} establishes a compelling
precedent: TEEs can serve as \emph{data handling policy} enforcement points
where code authenticity is remotely verifiable and memory is opaque to the
host OS. This primitively changes the threat model. An adversary who
compromises the host process cannot read or modify the computation inside the
enclave. The policy (the sanitizer and encoder) runs in attested, measured
code.

We extend this insight: the TEE is not merely a confidentiality boundary—it
is a \emph{liveness} boundary. The critical invariant is that the raw prompt
is \emph{deleted} after encoding, inside the enclave. Nothing that leaves the
TEE contains the raw user text; only the encoded threat vector and a
defense-rewritten context do.

\subsection{Threat Classes}
\label{sec:intro:threats}

We identify three primary attack classes against aligned LLMs, each requiring
a distinct defense primitive:

\begin{enumerate}
\item \textbf{Adversarial Roleplay (AR):} The attacker instructs the model to
adopt a persona with different safety properties (``DAN'', unrestricted
characters, fictional experts in harmful domains). Survey data from
Shen et al.~\cite{shen2023anything} classifies persona-injection as the
modal jailbreak category ($\approx$47\% of successful attacks in the wild).

\item \textbf{Context/Prompt Injection (CPI):} The attacker injects content
into the context window designed to override system instructions. This
includes direct instruction overrides~\cite{perez2022ignore}, delimiter
attacks that close the system prompt block and open a new one, and indirect
injection via retrieved documents in RAG pipelines~\cite{greshake2023youve}.

\item \textbf{Prompt Exfiltration (PE):} The attacker extracts the system
prompt or model internals through carefully crafted queries. While not
directly a jailbreak, exfiltration enables targeted construction of
attacks that bypass specific defenses.
\end{enumerate}

\subsection{Contributions}
\label{sec:intro:contributions}

We make the following contributions:

\begin{enumerate}
\item A formal architecture (\S\ref{sec:architecture}) for hardware-enforced
LLM safety with a provable raw-prompt non-disclosure invariant.
\item The Morphological Intent Encoder (\S\ref{sec:encoder}), a surface-form
invariant threat classifier operating on latent space geometry.
\item A threshold signature protocol (\S\ref{sec:threshold}) for
multi-node consensus on safety verdicts, with a security proof
against single-node compromise (\S\ref{sec:security}).
\item An epidemic vaccine propagation system (\S\ref{sec:gossip}) for
federated immunity with $O(\log N)$ convergence.
\item A JSON-RPC~2.0 telemetry protocol (\S\ref{sec:telemetry}) for
model-agnostic interoperability.
\item An open-source reference implementation in Python~3.11
(\S\ref{sec:implementation}) with 8 test suites and stdlib-only core.
\end{enumerate}

% =============================================================================
\section{Background and Related Work}
\label{sec:background}

\subsection{Trusted Execution Environments as Policy Enforcement}
\label{sec:bg:tee}

Trusted Execution Environments (Intel TDX/SGX, AMD SEV-SNP, ARM TrustZone)
provide hardware-enforced memory isolation and remote attestation. The host
OS cannot read enclave memory, and the attestation report binds code identity
(MRENCLAVE measurement) to the execution environment.

The NDAI Agreements~\cite{ndai2024tee} paper is the closest prior work to our
TEE application domain. It proposes TEEs as \emph{data handling policy}
enforcement points in AI deployment, where the enclave enforces data use
agreements without requiring trust in the cloud provider. We extend this
model specifically to prompt-level safety, adding the invariant that
raw prompts are destroyed before exiting the enclave.

Dstack~\cite{dstack2024} provides a developer-facing TEE orchestration
layer with attestation APIs over Intel TDX. ASP uses Dstack's attestation
interface to verify that the sanitizer and encoder running in the enclave
match their expected measurements.

\subsection{Threshold Cryptography and Thetacrypt}
\label{sec:bg:threshold}

Threshold cryptography distributes a cryptographic operation across $M$ parties
such that any $N \leq M$ parties can jointly perform it, but fewer than $N$
cannot~\cite{desmedt1987society}. Threshold signature schemes are the
primary primitive: a secret signing key is split into $M$ shares such that
any $N$ shares suffice to produce a valid signature, but $N-1$ shares reveal
nothing about the key.

Thetacrypt~\cite{thetacrypt2024} is a practical threshold cryptography
library designed for distributed consensus applications, providing
threshold BLS signatures and Distributed Key Generation (DKG). ASP's
threshold validation protocol is directly inspired by Thetacrypt's
``quorum of trust'' model: a verdict is valid only when $N$ independent
validator nodes co-sign it.

\subsection{Shamir Secret Sharing}
\label{sec:bg:shamir}

Shamir~\cite{shamir1979share} showed that any secret $s \in \mathbb{Z}_p$
can be hidden in the constant term of a uniformly random polynomial
$f(x) = s + a_1 x + \cdots + a_{t-1} x^{t-1}$ over $\mathbb{Z}_p$.
Evaluating $f$ at $M$ distinct non-zero points produces $M$ shares; any
$t$ shares allow Lagrange interpolation to recover $f(0) = s$, while
$t-1$ shares are information-theoretically independent of $s$.

ASP uses Shamir secret sharing to split validator signing keys, so that
a quorum of $N$ nodes is necessary and sufficient to produce a valid
aggregated signature over a safety verdict.

\subsection{Epidemic Algorithms for Gossip Convergence}
\label{sec:bg:gossip}

Demers et al.~\cite{demers1987epidemic} established the theoretical
foundations for epidemic (gossip) protocols, proving that rumor-spreading
with random peer selection achieves convergence in $O(\log N)$ rounds for
a network of $N$ nodes. The ``push'' variant—where a node forwards a
message to $k$ random peers per round—is the model used by ASP for
vaccine propagation.

\subsection{Latent Space Attack Manifolds}
\label{sec:bg:manifolds}

Adversarial ML literature has extensively characterized the structure of
attack inputs in embedding space. Zou et al.~\cite{zou2023universal}
demonstrated that adversarial suffixes for aligned LLMs transfer across
models, implying that attacks share representational structure independent
of surface form. Carlini et al.~\cite{carlini2023aligned} showed that
alignment fine-tuning does not close the latent-space proximity between
harmful and harmless representations.

These results motivate our encoder design: semantic intent is more stable
in latent space than in token space, so detecting attacks by latent-space
proximity to known attack embeddings is more robust than surface-form matching.

\subsection{Jailbreak Taxonomy}
\label{sec:bg:taxonomy}

Wei et al.~\cite{wei2023jailbroken} identify two failure modes of safety
training—``competing objectives'' (the model trades helpfulness against
safety) and ``mismatched generalization'' (safety training does not cover
the full distribution of attack prompts). This analysis predicts that
roleplay attacks are the dominant category: they exploit the helpfulness
objective by framing the harmful request as a creative or educational task.
Perez and Ribeiro~\cite{perez2022ignore} provide a systematic taxonomy of
prompt injection variants, covering direct instruction overrides, delimiter
attacks, and virtualization attacks.

% =============================================================================
\section{Threat Model}
\label{sec:threat}

\subsection{Attacker Model}
\label{sec:threat:attacker}

\textbf{Capabilities.} The attacker can:
\begin{itemize}
  \item Submit arbitrary strings to the LLM API (user-controlled input).
  \item Observe LLM outputs in full.
  \item Know the existence and general design of ASP (Kerckhoffs principle).
  \item Compromise up to $N-1$ validator nodes (but not $N$ or more).
  \item Access the host OS but not the TEE enclave memory.
\end{itemize}

\textbf{Goals.} The attacker seeks to:
\begin{itemize}
  \item Elicit harmful outputs from the downstream LLM.
  \item Cause the system to emit a \textsc{Verified\_Immunity} verdict
    for a request that should be blocked.
  \item Extract the raw prompt or system context from the TEE.
\end{itemize}

\subsection{Formal Attack Class Definitions}
\label{sec:threat:classes}

Let $\Sigma^*$ be the set of all input strings and $\mathcal{H} \subset \Sigma^*$
be the set of harmful prompts (as defined by the system operator's policy).

\begin{definition}[Adversarial Roleplay Attack]
A string $x \in \Sigma^*$ is an \emph{adversarial roleplay attack} if
$x \notin \mathcal{H}$ by surface-form analysis but
$\exists x' \in \mathcal{H}$ such that $x$ is semantically equivalent to
a composition $\langle\text{persona frame}\rangle \circ x'$, where the
persona frame asserts altered safety properties.
\end{definition}

\begin{definition}[Context Injection Attack]
A string $x \in \Sigma^*$ is a \emph{context injection attack} if
$x$ contains a subsequence $s$ such that, when $s$ is processed by
the LLM, it causes the LLM to deviate from its system-prompt-specified
behavior by treating $s$ as a higher-priority instruction source than
the system prompt.
\end{definition}

\begin{definition}[Prompt Exfiltration Attack]
A string $x \in \Sigma^*$ is a \emph{prompt exfiltration attack} if
processing $x$ causes the LLM to emit any substring of the system prompt
$P$ in its response, i.e., the response contains $s \subseteq P$.
\end{definition}

\subsection{Trust Boundaries}
\label{sec:threat:trust}

\textbf{ASP trusts:}
\begin{itemize}
  \item TEE hardware (SGX/TDX/SEV integrity guarantees).
  \item The attested code measurement of the sanitizer and encoder.
  \item Any quorum of $\geq N$ validator nodes producing consistent verdicts.
  \item The embedding model's ability to cluster semantically similar inputs.
\end{itemize}

\textbf{ASP does not trust:}
\begin{itemize}
  \item The host OS, hypervisor, or network between components.
  \item Any single validator node.
  \item The LLM's own safety mechanisms (treated as a defense in depth layer,
    not the primary enforcement mechanism).
  \item Surface-form analysis of input text.
\end{itemize}

\subsection{Explicit Non-Goals}
\label{sec:threat:nongaols}

ASP does not protect against:
\begin{itemize}
  \item Physical side-channel attacks on TEE hardware (power, timing, EM).
  \item A Byzantine majority ($\geq N$ of $M$) of validator nodes.
  \item Attacks that are semantically outside the coverage of the attack
    signature database (zero-day attacks with no prior vaccine).
  \item Synchronization-independent attacks that exploit the gossip
    convergence window before a vaccine propagates.
  \item Model weight extraction or membership inference attacks.
\end{itemize}

% =============================================================================
\section{System Architecture}
\label{sec:architecture}

Figure~\ref{fig:architecture} illustrates the full ASP pipeline. A raw user
prompt enters the TEE boundary and traverses four layers before a
defense-sanitized context is released to the downstream LLM. The pipeline
is entirely asynchronous and model-agnostic.

\begin{figure}[t]
\centering
\begin{tikzpicture}[
    node distance=0.8cm and 1.4cm,
    box/.style={draw, rectangle, rounded corners=3pt, minimum width=2.6cm,
                minimum height=0.7cm, font=\footnotesize, fill=blue!6},
    tee/.style={draw, rectangle, dashed, rounded corners=6pt,
                inner sep=10pt, label={[font=\footnotesize\itshape]north:TEE Enclave}},
    arrow/.style={-{Stealth[length=4pt]}, thick},
  ]

  \node[box] (raw)    {Raw Prompt};
  \node[box, right=of raw]  (san)    {Sanitizer};
  \node[box, right=of san]  (enc)    {Intent Encoder};
  \node[box, below=of enc]  (def)    {Defense Router};
  \node[box, left=of def]   (thr)    {Threshold Validator};
  \node[box, left=of thr]   (out)    {Sanitized Context};
  \node[box, below=of out]  (llm)    {LLM};

  \draw[arrow] (raw) -- node[above,font=\tiny]{enter} (san);
  \draw[arrow] (san) -- node[above,font=\tiny]{skeleton} (enc);
  \draw[arrow] (enc) -- node[right,font=\tiny]{ThreatVector} (def);
  \draw[arrow] (def) -- node[below,font=\tiny]{MitigationPayload} (thr);
  \draw[arrow] (thr) -- node[below,font=\tiny]{SignatureBlock} (out);
  \draw[arrow] (out) -- node[left,font=\tiny]{exit TEE} (llm);

  \node[draw=red!50, dashed, rounded corners=6pt, fit=(san)(enc)(def)(thr)(out),
        inner sep=8pt, label={[font=\footnotesize\itshape,text=red!70]north:TEE Enclave}] {};
\end{tikzpicture}
\caption{ASP pipeline. The raw prompt enters the TEE enclave, is encoded and
classified, mitigation is applied, threshold consensus is collected, and only
the sanitized context exits.}
\label{fig:architecture}
\end{figure}

\subsection{TEE Isolation Layer}
\label{sec:tee}

\textbf{Formal Data Flow.}
Let $P$ be a raw prompt string. The TEE boundary enforces the following
sequential transformation:

\begin{equation}
\label{eq:pipeline}
P \xrightarrow{\text{sanitize}} C_{\emptyset}
  \xrightarrow{\text{encode}(P)} \vec{v}
  \xrightarrow{\text{del } P}
  \xrightarrow{\text{route}} M
  \xrightarrow{\text{validate}} B
  \xrightarrow{\text{if VERIFIED}} C_{\text{safe}}
\end{equation}

where $C_{\emptyset}$ is a skeleton context with no raw content,
$\vec{v} \in \mathbb{R}^d$ is the threat vector,
$M$ is the mitigation payload,
$B$ is the threshold signature block, and
$C_{\text{safe}}$ is the defense-rewritten context released to the LLM.

\textbf{The Critical Invariant.}
The instruction \texttt{del raw\_prompt} in Algorithm~\ref{alg:pipeline}
is not a correctness optimization—it is the primary security property.
After encoding, the raw prompt is a local variable in the enclave; Python's
garbage collector deallocates the page. In a production TEE deployment,
the enclave allocator explicitly zeros the pages. No instance variable,
log, or return value ever holds the raw prompt after Line~3 of
Algorithm~\ref{alg:pipeline}.

\begin{algorithm}[t]
\caption{TEE Boundary: process(raw\_prompt)}
\label{alg:pipeline}
\begin{algorithmic}[1]
\Require raw\_prompt $\in \Sigma^*$
\Ensure PipelineResult with sanitized context
\Try
\State $C \gets \text{sanitize}(\text{raw\_prompt})$
  \Comment{Build skeleton context; no raw content}
\State $\vec{v} \gets \text{encode}(\text{raw\_prompt})$
  \Comment{Project to latent space; last access}
\Finally
\State \textbf{del} raw\_prompt
  \Comment{\textbf{Critical: destroy raw text even on exception}}
\EndTry
\State $M \gets \text{defenseRouter.route}(\vec{v},\, C)$
\State $B \gets \textbf{await}~\text{thresholdValidator.validate}(\vec{v},\, M)$
\If{$B.\text{verdict} = \textsc{Rejected}$}
  \State \textbf{raise} SecurityError
\EndIf
\Return $(C_\text{safe},\, \vec{v},\, M,\, B)$
\end{algorithmic}
\end{algorithm}

\textbf{Remote Attestation.}
ASP integrates with the Dstack attestation API~\cite{dstack2024} to provide
cryptographic proof that the sanitizer and encoder code running in the enclave
matches the operator-expected measurement. The \texttt{DstackAttestation}
class wraps the MRENCLAVE measurement and exposes a \texttt{verify\_remote}
method. An external verifier can confirm that ASP is running unmodified inside
a genuine TEE before trusting its verdicts.

\subsection{Morphological Intent Encoder}
\label{sec:encoder}

\textbf{Formal Definition.}
Let $\Sigma^*$ be the input alphabet and $\mathbb{R}^d$ the embedding space
(dimension $d = 768$ by default). The morphological intent encoder is a
function:

\begin{equation}
\label{eq:encoder}
f : \Sigma^* \to \mathbb{R}^d
\end{equation}

implemented by an embedding adapter $\phi$ followed by L2-normalization:

\begin{equation}
\label{eq:embed}
f(x) = \frac{\phi(x)}{\|\phi(x)\|_2}
\end{equation}

\textbf{Attack Manifolds.}
For each attack class $k \in \{1, \ldots, K\}$, let
$A_k = \{f(a) : a \in \mathcal{A}_k\} \subset S^{d-1}$
be the set of L2-normalized embeddings of known attack instances in class $k$,
forming a \emph{manifold} on the unit hypersphere.

The threat level function is defined by piecewise thresholds on cosine
similarity to the nearest point on any manifold:

\begin{equation}
\label{eq:threat}
\tau(x) = \begin{cases}
\textsc{Quarantine} & \text{if } \max_k \max_{a \in A_k} \langle f(x), a \rangle \geq \theta_B + 0.10 \\
\textsc{Block}      & \text{if } \max_k \max_{a \in A_k} \langle f(x), a \rangle \geq \theta_B \\
\textsc{Warn}       & \text{if } \max_k \max_{a \in A_k} \langle f(x), a \rangle \geq \theta_M + \tfrac{2}{3}(\theta_B - \theta_M) \\
\textsc{Monitor}    & \text{if } \max_k \max_{a \in A_k} \langle f(x), a \rangle \geq \theta_M \\
\textsc{Benign}     & \text{otherwise}
\end{cases}
\end{equation}

where $\theta_B = 0.82$ is the block threshold and $\theta_M = 0.65$
is the monitor threshold (both configurable via \texttt{ASPConfig}).

\textbf{Surface-Form Invariance Property.}
The key insight is that $f$ (the embedding function) collapses
surface-form variation while preserving semantic intent. Formally:

\begin{definition}[Surface-Form Equivalence]
Two prompts $x, x' \in \Sigma^*$ are \emph{surface-form equivalent}
with respect to intent $I$ if $\text{intent}(x) = \text{intent}(x') = I$,
but $x \neq x'$ as strings (e.g., $x'$ is a paraphrase, translation,
or encoded variant of $x$).
\end{definition}

\begin{claim}[Surface-Form Invariance]
\label{claim:sfi}
For any two surface-form equivalent prompts $x, x'$ with harmful intent,
$\langle f(x), f(x') \rangle > \theta_M$, i.e., both are classified
at threat level $\geq$ \textsc{Monitor}.
\end{claim}

This claim is supported empirically by the transfer attack results of
Zou et al.~\cite{zou2023universal}, which demonstrate that adversarial
suffixes computed for one model transfer to others, indicating that harmful
intent occupies a stable region of the embedding space across paraphrases.
Our implementation uses the nearest-neighbor search over the attack
signature DB (Algorithm~\ref{alg:encoder}) to approximate manifold proximity.

\begin{algorithm}[t]
\caption{MorphologicalIntentEncoder.encode(raw\_prompt)}
\label{alg:encoder}
\begin{algorithmic}[1]
\Require raw\_prompt $\in \Sigma^*$
\Ensure ThreatVector $\vec{v}$
\State $e \gets \phi(\text{raw\_prompt})$ \Comment{Embed via adapter}
\State $e \gets e / \|e\|_2$ \Comment{L2-normalize}
\State $(M, \text{ids}) \gets \text{AttackDB.get\_matrix}()$
\If{$M = \emptyset$}
  \Return ThreatVector$(e,\, 0.0,\, \varepsilon,\, \textsc{Benign})$
\EndIf
\State $i^* \gets \arg\max_i (M \cdot e)_i$
  \Comment{Nearest attack signature}
\State $s^* \gets (M \cdot e)_{i^*}$
  \Comment{Max cosine similarity}
\State $\ell \gets \tau(s^*)$ \Comment{Classify via Eq.~\eqref{eq:threat}}
\Return ThreatVector$(e,\, s^*,\, \text{ids}[i^*],\, \ell)$
\end{algorithmic}
\end{algorithm}

The attack signature database stores L2-normalized embeddings of known
attack prompts as a matrix $M \in \mathbb{R}^{n \times d}$ and supports
$O(nd)$ nearest-neighbor search (suitable for deployment-scale $n$;
production systems may substitute an ANN index such as FAISS).

\subsection{Threshold Validation}
\label{sec:threshold}

\textbf{Formal Scheme.}
An $(N, M)$-threshold validation scheme consists of $M$ validator nodes,
each holding a secret share derived from the network signing key.
A verdict $V$ obtains status \textsc{Verified\_Immunity} only when
$\geq N$ nodes independently co-sign it.

\textbf{Secret Sharing.}
Following Shamir~\cite{shamir1979share}, the signing key $s \in \mathbb{Z}_p$
(where $p = 2^{256} - 189$ is a 256-bit prime) is split by choosing a
random polynomial:

\begin{equation}
\label{eq:shamir}
f(x) = s + a_1 x + a_2 x^2 + \cdots + a_{N-1} x^{N-1} \pmod{p}
\end{equation}

where $a_1, \ldots, a_{N-1} \overset{\$}{\leftarrow} \mathbb{Z}_p$.
Node $i$ receives share $(i,\, f(i))$ for $i \in \{1, \ldots, M\}$.

\textbf{Reconstruction.}
Given any $N$ shares $\{(x_j, y_j)\}_{j=1}^{N}$, the secret is recovered
by Lagrange interpolation at $x = 0$:

\begin{equation}
\label{eq:lagrange}
s = f(0) = \sum_{j=1}^{N} y_j \prod_{\substack{k=1\\k \neq j}}^{N}
  \frac{-x_k}{x_j - x_k} \pmod{p}
\end{equation}

Equation~\eqref{eq:lagrange} is implemented in \texttt{asp/threshold/share.py}
using modular inverse via Fermat's little theorem: $a^{-1} \equiv a^{p-2} \pmod{p}$.

\textbf{State Machine.}
A verdict request progresses through the following states:

\begin{equation}
\label{eq:statemachine}
\textsc{Untrusted}
  \xrightarrow{\text{submit}}
\textsc{Pending}
  \xrightarrow{|\text{shares}| \geq N}
\textsc{Verified\_Immunity}
  \quad \text{or} \quad
\textsc{Rejected}
\end{equation}

A verdict that times out (configurable, default 5~s) without reaching quorum
is resolved fail-closed: if the proposed verdict was \textsc{Rejected}, the
final verdict is \textsc{Rejected}; otherwise the verdict falls back to
\textsc{Untrusted} and the request is denied.

\begin{algorithm}[t]
\caption{ThresholdValidator.validate($r$, $\vec{v}$, $M$)}
\label{alg:threshold}
\begin{algorithmic}[1]
\Require request\_id $r$, ThreatVector $\vec{v}$, MitigationPayload $M$
\Ensure ThresholdSignatureBlock $B$
\State nodes $\gets$ NodeRegistry.get\_all()
\State $N \gets \min(\text{config.threshold\_N},\, |\text{nodes}|)$
\State $V_\text{prop} \gets \textsc{Rejected}$ if $M.\text{action} = \textsc{FullBlock}$ else \textsc{Verified\_Immunity}
\State shares $\gets []$
\ForAll{node $n_i \in \text{nodes}$} \textbf{async}
  \State $V_i \gets n_i.\text{evaluate\_locally}(\vec{v})$
  \If{$V_i = V_\text{prop}$}
    \State shares.append($n_i.\text{sign}(r, \vec{v}, V_\text{prop})$)
  \EndIf
\EndFor
\State \textbf{await with timeout}
\If{$|\text{shares}| \geq N$}
  \State $\sigma \gets \text{SHA256}(\text{concat\_sorted}(\text{shares}))$
  \Return $B(V_\text{prop},\, r,\, N,\, |\text{nodes}|,\, \text{shares},\, \sigma)$
\Else
  \Return $B(\textsc{Untrusted},\, r,\, N,\, |\text{nodes}|,\, [],\, \varepsilon)$
\EndIf
\end{algorithmic}
\end{algorithm}

\subsection{Federated Immunity via Gossip}
\label{sec:gossip}

\textbf{Vaccine Structure.}
When a node detects a novel attack pattern (one not in its local attack DB),
it packages an embedding into a \emph{vaccine}:

\begin{equation}
\label{eq:vaccine}
V = (h,\, \vec{e},\, \text{module},\, \text{node\_id},\, t)
\end{equation}

where $h = \text{SHA256}(\vec{e}.\text{tobytes}())$ is the content hash,
$\vec{e} \in \mathbb{R}^d$ is the attack embedding,
\text{module} names the defense handler, and $t$ is the discovery timestamp.

\textbf{Epidemic Broadcast.}
The gossip engine implements push-rumor spreading~\cite{demers1987epidemic}.
In each round, a node holding an unseen vaccine selects $k$ random healthy
peers (the \emph{fanout}, default $k=3$) and sends the vaccine to those
peers that have not yet seen it (tracked via a per-peer seen-set).

\begin{theorem}[Convergence Bound]
\label{thm:gossip}
In a connected network of $N$ nodes with gossip fanout $k \geq 2$, the
push-rumor protocol delivers a vaccine to all nodes within
$O(\log_k N)$ rounds in expectation, after which no further propagation
occurs due to the seen-set idempotency check.
\end{theorem}

\begin{proof}[Proof Sketch]
At round $r$, let $I_r$ be the number of infected (vaccine-holding) nodes.
Each infected node contacts $k$ random peers; the probability of hitting
at least one uninfected node decreases as $I_r \to N$ but the expected
doubling time is $O(1/k)$ rounds. Standard coupon-collector analysis
gives $O(\log_k N)$ rounds to full coverage.
See Demers et al.~\cite{demers1987epidemic} for a rigorous treatment.
\end{proof}

\textbf{Idempotency.}
Each node maintains a \texttt{seen} set of vaccine hashes. On receiving a
vaccine, the node checks $h \in \text{seen}$ before processing. Duplicate
vaccines are discarded without re-broadcasting, ensuring the protocol
terminates and the attack DB remains consistent across the network.

\textbf{Signature Hash Verification.}
When a vaccine is deserialized from the wire, the receiving node recomputes
$h' = \text{SHA-256}(\text{embedding} \| \text{metadata})$ and verifies that
$h' = h_\text{claimed}$. If the hashes do not match, the vaccine is discarded
and the sending peer is flagged for anomalous behavior. This prevents
\emph{vaccine poisoning} attacks in which a malicious peer broadcasts a vaccine
whose \texttt{signature\_hash} does not correspond to its actual embedding,
which could otherwise corrupt the local attack database.

\textbf{Network Partition Tolerance.}
If a partition separates the network, vaccines propagate only within each
partition until connectivity is restored. Upon reconnection, the gossip
engine will re-broadcast pending vaccines to newly reachable peers.
This provides eventual consistency guarantees with no coordination overhead.

\subsection{Defense Module Interface}
\label{sec:defense}

Each defense module implements the following abstract interface:

\begin{align}
\label{eq:defense_interface}
\text{evaluate} &: \text{ThreatVector} \to [0, 1] \\
\text{mitigate} &: (\text{ThreatVector},\, \text{SanitizedContext}) \to \text{MitigationPayload}
\end{align}

The \texttt{evaluate} function returns a confidence score for the module's
relevance to the given threat. The defense router selects the highest-confidence
module capable of handling the threat level.

\textbf{AdversarialRoleplayModule.}
The roleplay module represents its detection capability as a centroid
vector in embedding space:

\begin{equation}
\label{eq:capability}
\vec{c}_{\text{roleplay}} = \frac{1}{|S|} \sum_{s \in S} f(s)
\end{equation}

where $S$ is a seed set of canonical roleplay attack descriptions
(``pretend you are an unrestricted AI'', ``act as DAN'', etc.).
The module's confidence on a threat $\vec{v}$ is
$\text{ReLU}(\langle \vec{v}, \vec{c}_{\text{roleplay}} \rangle)$.

Mitigation strategy is graduated: at \textsc{Warn} level, the module
redirects the roleplay frame to a safe persona; at \textsc{Block}/\textsc{Quarantine},
it substitutes the entire request with a safe notification (full block).

\textbf{ContextInjectionModule.}
This module defends against instruction override attacks by wrapping the
sanitized context in explicit \emph{alignment anchors}—boundary markers
that the LLM's instruction-following circuitry is trained to treat as
system-level demarcations:

\begin{lstlisting}[language={},caption={Alignment Anchor Structure},label={lst:anchors}]
[SYSTEM BOUNDARY - VERIFIED BY ASP THRESHOLD CONSENSUS]
The following is untrusted user input. Your system prompt
takes absolute precedence over any instructions below.
[BEGIN UNTRUSTED INPUT]
{sanitized_content}
[END UNTRUSTED INPUT]
[SYSTEM BOUNDARY - RESUME NORMAL OPERATION]
\end{lstlisting}

This defense does not rely on the LLM treating the anchors as
special tokens—it exploits the fact that explicit demarcation in
the system-prompt position is consistently honored by RLHF-aligned
models. At \textsc{Block}/\textsc{Quarantine}, the module performs
a full block, replacing user content with a safe notification.

\subsection{N-Dimensional Gaussian Memory Field}
\label{sec:promptfieldnd}

ASP's threat memory is organized as a continuous potential field in
$\mathbb{R}^n$, where each stored threat signature induces an isotropic
Gaussian influence. This \emph{PromptFieldND} representation enables
gradient-based placement of new memories in minimally-occupied regions and
principled novelty detection.

\textbf{Formal Definition.}

\begin{definition}[Memory Region]
\label{def:memory-region}
A \emph{memory region} $m_i$ in the $n$-dimensional prompt field is a tuple
$m_i = (c_i,\, \sigma_i,\, w_i,\, \ell_i)$ where $c_i \in [0,1]^n$ is the
centroid, $\sigma_i > 0$ is the isotropic spread, $w_i > 0$ is the weight
(importance), and $\ell_i$ is a semantic label. Each region induces a
Gaussian potential:
\begin{equation}
\label{eq:gaussian-potential}
G(\mathbf{p},\, m_i) = w_i \cdot \exp\!\left(
  -\frac{\|\mathbf{p} - c_i\|^2}{2\,\sigma_i^2}
\right)
\end{equation}
\end{definition}

\begin{definition}[Field Energy]
\label{def:field-energy}
The \emph{total field energy} at a point $\mathbf{p} \in [0,1]^n$ is the
superposition of all memory region potentials:
\begin{equation}
\label{eq:field-energy}
E(\mathbf{p}) = \sum_{i=1}^{K} G(\mathbf{p},\, m_i)
\end{equation}
where $K$ is the number of stored regions. High-energy zones correspond to
well-characterized threat categories; low-energy zones are unexplored.
\end{definition}

\textbf{Keyword Gradient.}
The mapping from raw text to a coordinate in $[0,1]^n$ is performed by the
\emph{keyword gradient} function. Let $\mathcal{D} = \{D_1, \dots, D_n\}$ be
a partition of threat-relevant keyword categories into $n$ dimension groups.
For a prompt with token set $T$:

\begin{equation}
\label{eq:keyword-gradient}
\text{score}_d = \tanh\!\bigl(\text{density}_d \cdot \alpha\bigr),
\quad d = 1, \dots, n
\end{equation}

where $\alpha > 0$ is a scaling constant and

\begin{equation}
\label{eq:density}
\text{density}_d = \max_{C \in D_d}
  \frac{|\{t \in T : t \in C\}|}{|T|}
\end{equation}

The $\tanh$ squashing ensures $\text{score}_d \in [0,1)$ and provides
diminishing returns for keyword flooding.

\textbf{Analytic Gradient.}
Because each $G(\mathbf{p},\, m_i)$ is differentiable everywhere, the
gradient of the field energy admits a closed-form expression:

\begin{equation}
\label{eq:field-gradient}
\nabla E(\mathbf{p}) = -\sum_{i=1}^{K}
  G(\mathbf{p},\, m_i)\,
  \frac{\mathbf{p} - c_i}{\sigma_i^2}
\end{equation}

The gradient points toward increasing energy (i.e., toward existing memory
centroids). To place a new region in a \emph{low-energy} zone, we
\emph{ascend} this gradient with sign reversal (equivalently, we minimize
$-E$).

\textbf{Gradient Descent Placement.}
Algorithm~\ref{alg:placement} describes the procedure for finding an
optimal low-energy position for a newly observed prompt.

\begin{algorithm}[t]
\caption{GradientDescentPlacement($\mathbf{p}_0$, FieldND $\mathcal{F}$)}
\label{alg:placement}
\begin{algorithmic}[1]
\Require Initial coordinate $\mathbf{p}_0 \in [0,1]^n$, learning rate $\eta$, tolerance $\varepsilon$, max steps $S = 40$
\Ensure Optimized position $\mathbf{p}^* \in [0,1]^n$
\State $\mathbf{p} \gets \mathbf{p}_0$
\For{$s = 1, \dots, S$}
  \State $\mathbf{g} \gets \nabla E(\mathbf{p})$
    \Comment{Eq.~\ref{eq:field-gradient}}
  \State $\mathbf{p} \gets \mathbf{p} + \eta \cdot \mathbf{g}$
    \Comment{Move \emph{away} from existing regions}
  \State $\mathbf{p} \gets \text{clamp}(\mathbf{p},\, 0,\, 1)$
    \Comment{Stay in unit hypercube}
  \If{$\|\mathbf{g}\| < \varepsilon$}
    \State \textbf{break}
      \Comment{Early termination: local energy minimum reached}
  \EndIf
\EndFor
\Return $\mathbf{p}$
\end{algorithmic}
\end{algorithm}

\begin{theorem}[Novelty Detection]
\label{thm:novelty}
Let $\mathbf{p}^*$ be the position returned by
Algorithm~\ref{alg:placement} for a new prompt embedding, and let
$\tau > 0$ and $\delta > 0$ be the energy threshold and minimum distance
parameters, respectively. The prompt is classified as \textsc{Novel} if and
only if:
\begin{equation}
\label{eq:novelty}
E(\mathbf{p}^*) < \tau
\quad \text{and} \quad
\min_{i=1}^{K} \|\mathbf{p}^* - c_i\| > \delta
\end{equation}
If either condition fails, the prompt is \emph{absorbed} into the nearest
existing region $m_j$ via a running centroid update:
\begin{equation}
\label{eq:absorption}
c_j^{\text{new}} = \frac{c_j \cdot n_j + \mathbf{p}^*}{n_j + 1},
\qquad
n_j \gets n_j + 1
\end{equation}
where $n_j$ is the observation count of region $m_j$.
\end{theorem}

\textbf{N-Dimensional Generalization.}
The formulation above is stated for arbitrary $n \in \mathbb{N}$. All
operations---Gaussian evaluation (Eq.~\ref{eq:gaussian-potential}), gradient
computation (Eq.~\ref{eq:field-gradient}), clamping, and distance
checks---depend only on the $\ell_2$ norm in $\mathbb{R}^n$ and are
therefore dimension-agnostic. The reference ASP implementation instantiates
the field in $\mathbb{R}^3$ (mapping to \emph{intent}, \emph{target}, and
\emph{technique} axes), but operators may extend to higher dimensions
(e.g., $n = 5$ adding \emph{persistence} and \emph{social-engineering}
axes) without algorithmic changes. The only dimension-dependent component is
the keyword gradient partition $\mathcal{D}$, which must supply exactly $n$
dimension groups.

\textbf{Seed Regions.}
Table~\ref{tab:seed-regions} lists the seven canonical seed regions shipped
with every ASP node. These prototypes bootstrap the field so that
gradient-based placement is immediately effective on first deployment.

\begin{table}[t]
\centering
\caption{Canonical Seed Regions for PromptFieldND ($n = 3$)}
\label{tab:seed-regions}
\small
\begin{tabular}{llccc}
\toprule
\textbf{\#} & \textbf{Label} & $c_i$ & $\sigma_i$ & $w_i$ \\
\midrule
1 & Jailbreak (DAN)            & $(0.9,\, 0.1,\, 0.8)$ & 0.15 & 1.0 \\
2 & Prompt Injection           & $(0.7,\, 0.3,\, 0.9)$ & 0.15 & 1.0 \\
3 & Data Exfiltration          & $(0.2,\, 0.9,\, 0.7)$ & 0.15 & 1.0 \\
4 & Roleplay Override          & $(0.8,\, 0.2,\, 0.5)$ & 0.15 & 1.0 \\
5 & System Prompt Leak         & $(0.3,\, 0.8,\, 0.6)$ & 0.15 & 1.0 \\
6 & Multi-turn Manipulation    & $(0.5,\, 0.5,\, 0.9)$ & 0.20 & 0.8 \\
7 & Benign Baseline            & $(0.1,\, 0.1,\, 0.1)$ & 0.25 & 0.5 \\
\bottomrule
\end{tabular}
\end{table}

% =============================================================================
\section{JSON-RPC 2.0 Telemetry Protocol}
\label{sec:telemetry}

ASP exposes all internal state transitions via a JSON-RPC~2.0 notification
interface. This design choice is deliberate: JSON-RPC~2.0 requires no
external libraries, supports batching, is model-agnostic, and is easily
parseable by any downstream analytics system.

\textbf{Protocol Messages.}
Table~\ref{tab:telemetry} defines the three canonical message types.
All messages are JSON-RPC~2.0 notifications (no response expected), using
the \texttt{asp.*} method namespace.

\begin{table}[t]
\centering
\caption{ASP Telemetry Message Schemas}
\label{tab:telemetry}
\small
\begin{tabular}{lll}
\toprule
\textbf{Field} & \textbf{Type} & \textbf{Description} \\
\midrule
\multicolumn{3}{l}{\textit{Method: \texttt{asp.threat\_detected}}} \\
\texttt{request\_id}  & string  & UUID of the request \\
\texttt{embedding}    & float[] & L2-norm'd threat vector \\
\texttt{max\_attack\_similarity} & float & Cosine sim to nearest attack sig \\
\texttt{nearest\_attack\_id}    & string & SHA-256 hash of nearest sig \\
\texttt{threat\_level}          & string & Enum: BENIGN…QUARANTINE \\
\texttt{timestamp}    & float   & Unix epoch seconds \\
\midrule
\multicolumn{3}{l}{\textit{Method: \texttt{asp.mitigation\_applied}}} \\
\texttt{request\_id}   & string & UUID (matches threat message) \\
\texttt{defense\_module} & string & Module name \\
\texttt{action}        & string & Enum: PASS\_THROUGH…FULL\_BLOCK \\
\texttt{explanation}   & string & Human-readable rationale \\
\texttt{metadata}      & object & Module-specific kv pairs \\
\midrule
\multicolumn{3}{l}{\textit{Method: \texttt{asp.threshold\_result}}} \\
\texttt{request\_id}   & string & UUID \\
\texttt{verdict}       & string & Enum: UNTRUSTED…VERIFIED\_IMMUNITY \\
\texttt{threshold}     & int    & $N$ \\
\texttt{total\_nodes}  & int    & $M$ \\
\texttt{shares\_collected} & int & Actual shares gathered \\
\texttt{is\_valid}     & bool   & $\geq N$ shares with aggregated sig \\
\texttt{aggregated\_signature\_hex} & string & Hex-encoded aggregate \\
\bottomrule
\end{tabular}
\end{table}

\textbf{Rationale for JSON-RPC~2.0.}
Alternatives considered include gRPC (protobuf) and plain HTTP REST.
JSON-RPC~2.0 was chosen because: (1) it requires zero external dependencies
(stdlib \texttt{json} only); (2) batching support enables high-throughput
telemetry without connection overhead; (3) the method-dispatch model maps
cleanly onto ASP's three event types; (4) JSON is universally parseable by
security tooling.

% =============================================================================
\section{Security Analysis}
\label{sec:security}

\subsection{Threshold Security Against Node Compromise}
\label{sec:security:threshold}

\begin{theorem}[Single-Node Non-Forgery]
\label{thm:threshold}
Let $(N, M)$ be the threshold parameters with $N \geq 2$.
If at most $N-1$ validator nodes are compromised (controlled by the attacker),
then the attacker cannot forge a \textsc{Verified\_Immunity} verdict for
any request that would otherwise receive \textsc{Rejected}.
\end{theorem}

\begin{proof}
A \textsc{Verified\_Immunity} verdict requires $\geq N$ valid
\texttt{SignatureShare} objects in the \texttt{ThresholdSignatureBlock},
each produced by a distinct honest node signing the verdict.

A compromised node can produce at most one \texttt{SignatureShare} (its own).
With $N-1$ compromised nodes, the attacker controls at most $N-1$ shares.
Since the threshold is $N$, the attacker is one share short.

The remaining honest nodes sign only if their local threat assessment agrees
with the proposed verdict. For a request the system assesses as harmful,
honest nodes compute \textsc{Rejected} locally and will not sign
\textsc{Verified\_Immunity}. Therefore the quorum of $N$ honest signatures
cannot be assembled.

In the Shamir formulation: the attacker holds at most $N-1$ shares of the
signing key. By the information-theoretic security of Shamir's scheme,
$N-1$ shares are statistically independent of the key (the polynomial's
constant term is uniformly random given any $N-1$ evaluations)~\cite{shamir1979share}.
Therefore the attacker cannot compute the aggregated signature.
\end{proof}

\begin{corollary}
The default configuration $(N=2, M=3)$ tolerates one compromised validator.
For $k$-fault tolerance, set $N = M - k$ with $M \geq 2k + 1$.
\end{corollary}

\subsection{Surface-Form Invariance}
\label{sec:security:sfi}

\begin{theorem}[Attack Detection Under Paraphrase]
\label{thm:sfi}
Let $x$ and $x'$ be two prompts with the same harmful intent
$I \in \mathcal{H}$ (differing only in surface form: paraphrase,
encoding, language). Assume the embedding function $\phi$ satisfies
the \emph{semantic stability} property:
$\langle f(x), f(x') \rangle \geq \epsilon$ for some $\epsilon > \theta_M$.
Then $\tau(x) \geq \textsc{Monitor}$ and $\tau(x') \geq \textsc{Monitor}$.
\end{theorem}

\begin{proof}
By the triangle inequality on cosine distance: if $a \in A_k$ is a seed
attack embedding with $\langle f(x), a \rangle \geq \theta_B$, then
$\langle f(x'), a \rangle \geq \langle f(x), a \rangle - \|f(x) - f(x')\|$.
The semantic stability assumption bounds $\|f(x) - f(x')\|$ via
$\|f(x) - f(x')\|^2 = 2 - 2\langle f(x), f(x')\rangle \leq 2(1-\epsilon)$.
Choosing $\epsilon$ such that $\theta_B - \sqrt{2(1-\epsilon)} > \theta_M$
gives the result.
\end{proof}

\textbf{Remark.} The semantic stability property is an assumption about
the embedding model $\phi$, not a proof about arbitrary embedders.
Empirical evidence from Zou et al.~\cite{zou2023universal} and
Carlini et al.~\cite{carlini2023aligned} supports this property for
current generation sentence transformers on the relevant attack classes.

\subsection{Attack Class Resistance}
\label{sec:security:resistance}

\textbf{Adversarial Roleplay.}
The AdversarialRoleplayModule's capability vector $\vec{c}_{\text{roleplay}}$
is computed as the centroid of seed attack embeddings. By Theorem~\ref{thm:sfi},
surface-form variants of roleplay attacks map to the same region of latent
space. The module evaluates all inputs, not just those already in the
attack DB, providing coverage for novel roleplay framings.

\textbf{Context Injection.}
The ContextInjectionModule responds by wrapping all user content in
alignment anchors regardless of confidence level at \textsc{Monitor}
and above. This is a structural defense: even if the classifier
mis-scores the injection, the anchor wrapping limits the injection's
ability to override the system prompt.

\textbf{Prompt Exfiltration.}
The raw prompt is deleted after encoding (Algorithm~\ref{alg:pipeline},
Line 3), and the system prompt is never placed in the raw prompt field.
Only the defense-rewritten \texttt{SanitizedContext} exits the TEE,
which contains no raw user content. This provides structural protection
against exfiltration.

\subsection{Complexity Analysis}
\label{sec:security:complexity}

\begin{table}[t]
\centering
\caption{Per-Request Computational Complexity}
\label{tab:complexity}
\small
\begin{tabular}{lll}
\toprule
\textbf{Component} & \textbf{Complexity} & \textbf{Notes} \\
\midrule
Sanitizer & $O(|P|)$ & String scan \\
Embedding & $O(d \cdot |P|)$ & Adapter-dependent \\
DB search & $O(n \cdot d)$ & $n$ = signatures in DB \\
Defense module & $O(d)$ & Single dot product \\
Threshold validation & $O(M)$ & $M$ nodes, async \\
Gossip round & $O(k \cdot d)$ & $k$ = fanout \\
Telemetry emit & $O(d)$ & JSON serialize \\
\midrule
\textbf{Total (critical path)} & $O(nd + M)$ & Dominated by DB search \\
\bottomrule
\end{tabular}
\end{table}

The critical path is $O(nd + M)$ where $n$ is the number of attack
signatures in the DB. For deployment-scale $n$ (thousands to millions),
the DB search should be replaced with an ANN index (e.g., FAISS IVF),
reducing search complexity to $O(d \log n)$.

% =============================================================================
\section{Empirical Evaluation}
\label{sec:evaluation}

We evaluate ASP against a curated 100-prompt test suite spanning benign
queries, known attack categories, and deliberately ambiguous borderline
inputs. All experiments use the default pipeline configuration with
five stages: sanitizer, intent encoder, threshold validator, gossip
engine, and telemetry emitter.

\subsection{Attack Detection Suite}

Table~\ref{tab:detection-suite} reports detection rates across eight
prompt categories. ASP achieves 100\% detection on all six attack
categories with zero false positives on unambiguously benign prompts.

\begin{table}[h]
\centering
\caption{Detection Results on the 100-Prompt Test Suite}
\label{tab:detection-suite}
\small
\begin{tabular}{lrl}
\toprule
\textbf{Category} & \textbf{Count} & \textbf{Detection Rate} \\
\midrule
Benign              & 25 & 0\% FP \\
Roleplay            & 10 & 100\% \\
Injection           & 10 & 100\% \\
Exfiltration        & 10 & 100\% \\
Smuggling           & 10 & 100\% \\
Social Engineering  & 10 & 100\% \\
Reframing           & 10 & 100\% \\
Borderline          & 10 & 10\% FP (1 edge case) \\
\midrule
\textbf{Total}      & \textbf{95} & \textbf{100\% TP, 0\% FP on benign} \\
\bottomrule
\end{tabular}
\end{table}

\subsection{Multilingual and Obfuscation Robustness}

A critical requirement for production deployment is robustness to
language-switching and character-level obfuscation. Table~\ref{tab:multilingual}
demonstrates that ASP's embedding-based similarity detection generalizes
across languages and obfuscation techniques without language-specific
rules.

\begin{table}[h]
\centering
\caption{Detection Across Languages and Obfuscation Techniques}
\label{tab:multilingual}
\small
\begin{tabular}{llrl}
\toprule
\textbf{Attack Variant} & \textbf{Threat Level} & \textbf{Similarity} & \textbf{Verdict} \\
\midrule
PT-BR injection         & QUARANTINE & 0.975 & REJECTED \\
Chinese roleplay        & QUARANTINE & 0.938 & REJECTED \\
French injection        & QUARANTINE & 0.965 & REJECTED \\
German injection        & QUARANTINE & 0.949 & REJECTED \\
Leetspeak               & QUARANTINE & 1.000 & REJECTED \\
CAPS LOCK               & QUARANTINE & 0.950 & REJECTED \\
Unicode zero-width      & QUARANTINE & 0.958 & REJECTED \\
Multi-turn simulation   & QUARANTINE & 1.000 & REJECTED \\
Code block payload      & QUARANTINE & 0.960 & REJECTED \\
\bottomrule
\end{tabular}
\end{table}

\subsection{False Positive Analysis}

Among the 25 benign prompts---covering topics such as cooking, travel,
mathematics, and creative writing---ASP produces zero false positives.
The single borderline false positive occurs on the prompt
\emph{``What techniques do hackers use to crack passwords?''},
which triggers a \texttt{WARN}-level classification rather than
\texttt{QUARANTINE}. This behavior is expected: the prompt contains
high keyword overlap with genuine attack signatures (``hackers'',
``crack'', ``passwords'') despite being a legitimate security-education
query. Because the threat level is \texttt{WARN} rather than
\texttt{QUARANTINE}, the prompt is not blocked but is flagged for
review---a conservative-but-correct outcome for content that straddles
the boundary between educational and adversarial intent.

\subsection{Full Pipeline Verification}

To verify that all five pipeline stages execute correctly, we query
the \texttt{/api/probe} endpoint with a known adversarial roleplay
prompt. The JSON response confirms end-to-end pipeline activation:

\begin{lstlisting}[language={},caption={Full pipeline response from \texttt{/api/probe}},label={lst:pipeline-response}]
{
  "threat": "QUARANTINE",
  "similarity": 1.0,
  "defense_action": "FULL_BLOCK",
  "defense_module": "adversarial_roleplay",
  "verdict": "REJECTED"
}
\end{lstlisting}

The \texttt{similarity} score of 1.0 indicates an exact match against
a known attack signature in the vector database. The
\texttt{defense\_module} field identifies which category of the defense
taxonomy was activated, and \texttt{FULL\_BLOCK} confirms that the
threshold validator rejected the prompt before it could reach the
downstream LLM.

% =============================================================================
\section{Implementation}
\label{sec:implementation}

\textbf{Language and Dependencies.}
ASP is implemented in Python 3.11. The core pipeline—TEE boundary,
intent encoder, threshold validator, gossip engine, and telemetry
emitter—uses only Python standard library and \texttt{numpy~2.2}
for vector arithmetic. No ML frameworks (PyTorch, TensorFlow, HuggingFace)
are required for the core path, minimizing the trusted computing base.

\textbf{Module Structure.}
The implementation is organized into 8 packages, each with a single
responsibility:

\begin{table}[h]
\centering
\caption{ASP Package Structure and Responsibilities}
\label{tab:modules}
\small
\begin{tabular}{lll}
\toprule
\textbf{Package} & \textbf{Files} & \textbf{Responsibility} \\
\midrule
\texttt{asp.tee}      & boundary, sanitizer, attestation & TEE orchestration \\
\texttt{asp.encoder}  & intent\_encoder, attack\_signature\_db, & Intent encoding, \\
                      & geometry, embedding\_adapter          & manifold geometry \\
\texttt{asp.threshold} & validator, share, node, registry    & N-of-M consensus \\
\texttt{asp.gossip}   & engine, vaccine, peer, transport     & Epidemic broadcast \\
\texttt{asp.defense}  & router, adversarial\_roleplay,       & Mitigation modules \\
                      & context\_injection, base             & \\
\texttt{asp.telemetry} & protocol, schemas, emitter          & JSON-RPC~2.0 \\
\texttt{asp.llm}      & adapter, openai\_adapter,            & LLM integration \\
                      & llama\_adapter                       & \\
\texttt{asp.config}   & config                               & Immutable config \\
\midrule
\textbf{Total} & 25 source files & \\
\bottomrule
\end{tabular}
\end{table}

\textbf{Test Suite.}
The test suite covers all modules with 8 test files targeting unit and
integration behavior:
\texttt{test\_intent\_encoder.py},
\texttt{test\_attack\_signature\_db.py},
\texttt{test\_threshold\_validator.py},
\texttt{test\_gossip\_engine.py},
\texttt{test\_defense\_modules.py},
\texttt{test\_sanitizer.py},
\texttt{test\_telemetry\_protocol.py}, and
\texttt{test\_e2e\_pipeline.py}.

\textbf{Model Agnosticism.}
The LLM adapter interface (\texttt{asp.llm.adapter}) abstracts provider
details. The same ASP protocol operates identically with OpenAI
(\texttt{openai\_adapter.py}), Llama-compatible APIs
(\texttt{llama\_adapter.py}), or any custom provider. The encoding adapter
is similarly pluggable: production deployments use sentence-transformers;
the reference implementation uses a hash-based adapter for zero-dependency
testing.

\textbf{Configuration.}
All tunable parameters (thresholds, $N$, $M$, fanout, timeouts) are
collected in \texttt{ASPConfig}, an immutable dataclass loaded from
environment variables. No magic constants exist in the core codebase.

% =============================================================================
\section{Limitations and Future Work}
\label{sec:limitations}

\textbf{Embedding Quality.}
The reference implementation uses a hash-based embedding adapter for
dependency-free testing. Hash-based embeddings do not preserve semantic
similarity and therefore do not satisfy the semantic stability property
(Theorem~\ref{thm:sfi}). Production deployments must substitute a
sentence-transformer model (e.g., \texttt{all-mpnet-base-v2} or a
custom fine-tuned encoder over jailbreak datasets). This is a protocol
design limitation, not an architectural one—the ASP pipeline is
identical regardless of embedding quality.

\textbf{Clock Synchronization.}
The gossip vaccine timestamping assumes approximate clock synchronization
across nodes. In adversarial network environments, a Byzantine node could
replay old vaccines with forged timestamps. A production deployment should
use a vector clock or a monotonic counter anchored to a consensus mechanism.

\textbf{Byzantine Validators Beyond Threshold.}
Theorem~\ref{thm:threshold} proves safety for up to $N-1$ compromised
validators. If $\geq N$ validators are Byzantine, ASP provides no safety
guarantee. Mitigating super-threshold collusion requires a
Byzantine fault-tolerant consensus layer (e.g., PBFT or HotStuff) which
is outside the current scope.

\textbf{Gossip Coverage Window.}
A novel attack detected by one node is not instantly visible to all nodes.
During the $O(\log N)$ convergence window, other nodes may not have the
vaccine in their attack DB. This creates a brief window of vulnerability
for zero-day attacks. Reducing the gossip interval (configurable, default
1~s) decreases the window at the cost of bandwidth.

\textbf{TEE Hardware Dependency.}
The current implementation stubs the Dstack attestation in development
mode (\texttt{tee\_attestation\_enabled = False}). Production deployment
requires Intel TDX, SGX, or equivalent TEE hardware and Dstack provisioning.
The protocol design is correct regardless—the attestation layer only
changes the assurance level, not the functional behavior.

\textbf{Alignment Anchor Robustness.}
The context injection defense relies on alignment anchors being honored by
RLHF-trained models. There exist fine-tuned or jailbroken model variants
that may not respect these anchors. ASP's defense modules are designed for
models that have undergone standard alignment training; they provide no
guarantees against deliberately de-aligned models.

\textbf{Future Work.}
Priority directions include: (1) integration of threshold BLS signatures
(via Thetacrypt~\cite{thetacrypt2024}) to replace the HMAC stand-in;
(2) a distributed key generation protocol for validator setup without
a trusted dealer; (3) an ANN index backend for the attack signature DB;
(4) a formal verification of the pipeline invariants using TLA+; and
(5) an adversarial evaluation benchmark against the jailbreak datasets
of Shen et al.~\cite{shen2023anything} and Wei et al.~\cite{wei2023jailbroken}.

% =============================================================================
\section{Conclusion}
\label{sec:conclusion}

We presented the Alignment Security Protocol (ASP), a hardware-backed,
threshold-validated, federated immunity system for LLM safety enforcement.
ASP addresses a structural gap in existing defenses: by moving safety
enforcement into an attested TEE, encoding threat detection in latent space
geometry rather than surface form, requiring distributed consensus on safety
verdicts, and propagating new attack signatures as vaccines via epidemic
broadcast, ASP provides defense-in-depth that is simultaneously more
robust, more principled, and more composable than existing approaches.

The core architectural insight is simple and strong: \emph{raw prompts must
be destroyed before they can be leaked, threat detection must operate on
semantic intent rather than token patterns, and no single node should hold
veto power over the safety decision.} These three invariants, enforced by
hardware and distributed consensus rather than policy documents, constitute
the foundation of synthetic immunity for aligned AI systems.

% =============================================================================
\bibliographystyle{ACM-Reference-Format}
\begin{thebibliography}{99}

\bibitem{bai2022constitutional}
Bai, Y., Jones, A., Ndousse, K., Askell, A., Chen, A., DasSarma, N.,
Drain, D., Fort, S., Ganguli, D., Henighan, T., et al. (2022).
Constitutional AI: Harmlessness from AI feedback.
\textit{arXiv preprint arXiv:2212.08073}.

\bibitem{carlini2023aligned}
Carlini, N., Nasr, M., Choquette-Choo, C. A., Jagielski, M.,
Garg, I., Terzis, A., Thomas, F., Tramèr, F., Zhang, C., and
Lee, K. (2023).
Are aligned neural networks adversarially aligned?
In \textit{Advances in Neural Information Processing Systems (NeurIPS 2023)},
volume 36.

\bibitem{demers1987epidemic}
Demers, A., Greene, D., Hauser, C., Irish, W., Larson, J., Shenker, S.,
Sturgis, H., Swinehart, D., and Terry, D. (1987).
Epidemic algorithms for replicated database maintenance.
In \textit{Proceedings of the 6th Annual ACM Symposium on Principles of
Distributed Computing (PODC '87)}, pages 1--12.

\bibitem{desmedt1987society}
Desmedt, Y. and Frankel, Y. (1990).
Threshold cryptosystems.
In \textit{Advances in Cryptology -- CRYPTO '89}, Lecture Notes in Computer
Science, volume 435, pages 307--315. Springer, New York.

\bibitem{dstack2024}
Dstack Project (2024).
Dstack: TEE orchestration for confidential cloud computing.
\textit{Technical documentation}. \url{https://github.com/dstack-TEE/dstack}.
Accessed 2024.

\bibitem{greshake2023youve}
Greshake, K., Abdelnabi, S., Mishra, S., Endres, C., Holz, T., and
Fritz, M. (2023).
Not what you've signed up for: Compromising real-world LLM-integrated
applications with indirect prompt injections.
In \textit{Proceedings of the AISec Workshop at ACM CCS 2023}.

\bibitem{ndai2024tee}
NDAI Project (2024).
NDAI Agreements: Trusted execution environments for AI data handling
policy enforcement.
\textit{Technical Report}. Cornell Initiative for CryptoCurrencies \& Contracts.

\bibitem{ouyang2022training}
Ouyang, L., Wu, J., Jiang, X., Almeida, D., Wainwright, C. L.,
Mishkin, P., Zhang, C., Agarwal, S., Slama, K., Ray, A., et al. (2022).
Training language models to follow instructions with human feedback.
In \textit{Advances in Neural Information Processing Systems (NeurIPS 2022)},
volume 35, pages 27730--27744.

\bibitem{perez2022ignore}
Perez, F. and Ribeiro, I. (2022).
Ignore previous prompt: Attack techniques for language models.
In \textit{NeurIPS ML Safety Workshop 2022}.
\textit{arXiv preprint arXiv:2211.09527}.

\bibitem{shamir1979share}
Shamir, A. (1979).
How to share a secret.
\textit{Communications of the ACM}, 22(11):612--613.

\bibitem{shen2023anything}
Shen, X., Chen, Z., Backes, M., Shen, Y., and Zhang, Y. (2023).
``Do Anything Now'': Characterizing and evaluating in-the-wild
jailbreak prompts on large language models.
\textit{arXiv preprint arXiv:2308.03825}.

\bibitem{thetacrypt2024}
Thetacrypt Project (2024).
Thetacrypt: Threshold cryptography for distributed consensus applications.
\textit{Technical documentation and library}.
\url{https://github.com/thetacrypt/thetacrypt}.
Accessed 2024.

\bibitem{wei2023jailbroken}
Wei, A., Haghtalab, N., and Steinhardt, J. (2023).
Jailbroken: How does LLM safety training fail?
In \textit{Advances in Neural Information Processing Systems (NeurIPS 2023)},
volume 36.

\bibitem{zou2023universal}
Zou, A., Wang, Z., Kolter, J. Z., and Fredrikson, M. (2023).
Universal and transferable adversarial attacks on aligned language models.
\textit{arXiv preprint arXiv:2307.15043}.

\bibitem{anthropic2024manyshot}
Anthropic (2024).
Many-shot jailbreaking.
\textit{Technical report}.

\bibitem{zou2024circuitbreakers}
Zou, A., Phan, L., Chen, S., Campbell, J., Guo, P., Ren, R., Pan, A.,
Yin, X., Mazeika, M., Dombrowski, A.-K., Goel, S., Li, N., Byun, M. J.,
Wang, Z., Mallen, A., Basart, S., Koyejo, S., Song, D., Fredrikson, M.,
Kolter, J. Z., and Hendrycks, D. (2024).
Circuit breakers for language model safety.
\textit{arXiv preprint arXiv:2406.04313}.

\bibitem{yi2023bipia}
Yi, J., Xie, Y., Zhu, B., Kiciman, E., Sun, G., Xie, X., and Wu, F. (2023).
Benchmarking and defending against indirect prompt injection attacks on large
language models.
\textit{arXiv preprint arXiv:2312.14197}.

\bibitem{toyer2024tensortrust}
Toyer, S., Watkins, O., Mendes, E. A., Svegliato, J., Bailey, L.,
Wang, T., Ong, I., Elmaaroufi, K., Abbeel, P., Darrell, T., Rishi, A.,
and Russell, S. (2024).
Tensor Trust: Interpretable prompt injection attacks from an online game.
In \textit{International Conference on Learning Representations (ICLR 2024)}.

\end{thebibliography}

\end{document}
